\section{Open Challenges and Future Directions}
% Length: ~10%
% Focus: Look forward. Point out the unresolved issues in the newest agentic frameworks.

While recent advances improve LLM-based evaluation, current mitigation strategies cannot yet guarantee robust and fair assessment. To systematize these unresolved issues and propose actionable solutions, we organize the remaining challenges into technical limitations and operational constraints, followed by a progressive research roadmap.

\subsection{Technical and Epistemic Challenges}

A critical bottleneck restricting progress is the absence of standardized benchmarks specifically designed to test RAG evaluator reliability. As highlighted by \citet{Brehme_2025}, who catalog existing evaluation methods and datasets, the field lacks dedicated testing grounds incorporating adversarial noise. Consequently, building a universally reliable judge remains difficult; many frameworks overfit to specific tasks and fail to generalize across diverse domains~\cite{agent_judge, auto_eval_judge}.

Beyond generalization, evaluators suffer from high prompt sensitivity. \citet{comm_systems} demonstrate that judgments are easily skewed by superficial features like politeness or verbosity, causing inconsistent scoring. Interventions like multi-agent debate attempt to fix this but often provide a false sense of security. Gains in these systems frequently stem from ensembling effects or majority pressure rather than genuine reasoning, allowing persuasive but incorrect agents to force a false consensus~\cite{debate_limit}. Furthermore, agentic evaluation introduces severe tool-memory conflicts. When external tool outputs contradict the model's internal priors, particularly in STEM fields, standard prompt corrections fail to reconcile the discrepancy~\cite{tool_conflict}. 

\subsection{Societal and Operational Constraints}

LLM judges also suffer from persistent prior biases. Models inherently reflect the distribution of their pre-training data. While frameworks like MACD~\cite{macd} assign diverse personas to different agents, shifting prompts often redistributes bias rather than eliminating it, complicating fair cross-cultural evaluation.

Operationally, solving technical flaws through multi-step agentic pipelines drastically increases computational cost and latency~\cite{agent_survey}. Tool use and memory retention also introduce distinct security vulnerabilities, including prompt injection and sensitive data leakage. These factors severely restrict real-world deployment where response time and privacy are paramount.

\subsection{A Roadmap for Future Research}

Overcoming these limitations demands approaches that are not only technically robust but also viable for large-scale, real-world application. We outline a two-phase roadmap to guide future evaluation frameworks.

In the short term (1 to 2 years), the community must prioritize creating adversarial benchmarks specifically for RAG evaluators~\cite{Brehme_2025} to measure cascading tool failures accurately. Concurrently, efforts should focus on distilling complex multi-agent reasoning into lightweight, single-pass models to bypass current latency and cost deployment barriers.

In the medium term (3 to 5 years), research should target architectural interventions to resolve internal tool-memory conflicts natively, moving beyond surface-level prompt corrections. Additionally, the field must develop dynamically localized evaluators capable of adjusting to specific cultural frameworks, shifting away from monolithic models that are constrained by a homogeneous set of cultural assumptions.